\section{Gestión de Riesgos}

\subsection{Identificación de Riesgos Técnicos}
\begin{table}[H]
    \centering
    \begin{tabularx}{\textwidth}{l l X X}
        \toprule
        \textbf{ID} & \textbf{Riesgo} & \textbf{Impacto} & \textbf{Mitigación} \\
        \midrule
        R-01 & Pérdida de comunicación. & Crítico & Implementar modo "Heartbeat" y retorno a base automático. \\
        R-02 & Atascamiento en terreno. & Alto & Algoritmo de detección de deslizamiento y maniobras de escape. \\
        R-03 & Sobrecalentamiento. & Medio & Disipadores pasivos y corte de energía preventivo. \\
        \bottomrule
    \end{tabularx}
\end{table}

\section{Apéndices}
\subsection{Glosario Técnico Detallado}
\begin{itemize}
    \item \textbf{SLAM:} Método por el cual un robot puede construir un mapa y posicionarse en él simultáneamente.
    \item \textbf{Heurística:} Técnica para resolver problemas más rápidamente cuando los métodos clásicos son demasiado lentos.
\end{itemize}

\subsection{Diagramas Adicionales}
(Reservado para esquemas eléctricos, planos mecánicos y diagramas de flujo de software detallados).
